\section{Git-Historie}
\label{app:git}

Die Entwicklung wurde in getrennten lokalen Commits dokumentiert. Der folgende
Auszug wurde mit \code{git log --oneline --decorate --graph --all} erzeugt. Der
öffentliche GitLab-Stand wird erst nach ausdrücklicher Freigabe hochgeladen.

\begin{lstlisting}[language={}]
* 0b279b9 Fallstudie und reproduzierbare Abbildungen ergänzen
* 1e34819 GitLab Pages und Projektdokumentation ergänzen
* e12b666 Daten- und Interaktionstests dokumentieren
* b6c805e Visualisierungen und Zeitraumsteuerung überarbeiten
* 4376a21 Datenpipeline auf Mai 2025 erweitern
*   79f38a9 Hochschul-GitLab mit bestehender Projekthistorie verbinden
|\
| * ff86945 Initial commit
* 969d7f3 2. Zwischenstand überarbeitet
* 3882dc7 2. Abbildung detailreicher und Zwischenbericht angepasst
* 693c16b GitHub hinzugefügt
* 23c6a28 Elm code sowie Zwischenstand aktualisiert
* 2c8bada Ordnerstruktur und erste zwischenstände
\end{lstlisting}

Der Bericht gehört nach der inhaltlichen und visuellen Endprüfung zu einem
eigenen abschlie{\ss}enden Commit. Dessen Hash kann nicht Bestandteil des in
diesem Commit erzeugten PDFs sein. Der Merge-Commit \code{79f38a9} verbindet
die vorhandene Projekthistorie mit dem initialen Commit des öffentlichen
Hochschul-GitLab-Projekts, ohne frühere Arbeit zu überschreiben.

\section{KI-gestützter Arbeitsprozess}
\label{app:ki}

Dieser Anhang fasst die Zusammenarbeit mit OpenAI Codex zusammen. Er ist kein
wörtliches Gesprächsprotokoll, sondern beschreibt die schrittweise erteilten
Arbeitsaufträge und die daraus entstandenen Ergebnisse. Thema, fachliche
Prioritäten und Freigaben wurden vom Verfasser festgelegt. Codex unterstützte
bei Recherche, Implementierungsentwürfen, Prüfungen und Formulierungen; die
Zwischenergebnisse wurden jeweils betrachtet, besprochen und weiterentwickelt.

\subsection*{A: Themenwahl und Datenzugang}

\begin{enumerate}
  \item Zunächst wurden mehrere geeignete Energy-Charts-Themen gegenübergestellt.
    Danach wurde die Untersuchung von Deutschlands Importen, Exporten,
    Länderbeziehungen, Erzeugungsmix und Preisen als Leitfrage festgelegt.
  \item Anschlie{\ss}end wurden der Datenbankzugang und das bereitgestellte Schema
    \code{energycharts} geprüft. Es folgten begrenzte PostgREST-Abfragen. Dabei
    wurden Datenquelle, HTTP-Endpunkte und Zugangshinweise getrennt von den
    exportierten Analysedaten dokumentiert.
  \item Nach Rückfragen zu den Einheiten wurden die OpenAPI-Beschreibung von
    Energy-Charts und die tatsächlichen Wertebereiche geprüft. Erst danach
    wurden GW, MW, EUR/MWh und die Vorzeichenkonvention verbindlich festgelegt.
  \item Für erste Tests wurden kleine Ausschnitte geladen. Nach erfolgreicher
    Prüfung wurde die Pipeline auf den vollständigen Mai 2025 mit 744 Stunden
    und elf Partnerländern erweitert.
\end{enumerate}

\subsection*{B: Elm-Visualisierungen und Interaktion}

\begin{enumerate}
  \item Der erste Prototyp verband einen gerichteten Ländergraphen, eine
    Zeitreihe und eine Pixelmatrix. Vorläufige Python-Abbildungen wurden danach
    vollständig durch SVG-Ansichten aus Elm ersetzt; Python blieb nur für
    Datenexport und Validierung zuständig.
  \item Die Bedienung wurde gemeinsam schrittweise geprüft: Land im Netzwerk
    wählen, Stunde in der Zeitreihe setzen und Land plus Stunde über eine
    Matrixzelle auswählen. Unabhängige Anwendungsfälle wurden deutlicher
    voneinander getrennt, damit Auswahlzustände nicht als feste Klickfolge
    missverstanden werden.
  \item Bei der Besprechung des Zeitdiagramms fiel auf, dass absolute Werte und
    Skalen fehlten. Daraufhin kamen eine beschriftete GW-Y-Achse, eine eigene
    symmetrische Flussskala, absolute Detailwerte und relative Anteile hinzu.
  \item Die Professorenhinweise wurden anschlie{\ss}end als Änderungsaufgaben
    umgesetzt: kleinere Pfeile, stärkere Maximallinien, mehrere Zeitfenster,
    globale Matrixfarben, eine Nullwertfarbe ohne wei{\ss}e Lücken und sichtbarere
    Auswahlmarkierungen.
  \item Nach jedem grö{\ss}eren Schritt wurden Elm-Build, Datensatzregeln und die
    relevanten Interaktionen erneut geprüft. Die finalen Abbildungen entstanden
    reproduzierbar aus festgelegten Zuständen der laufenden Elm-Anwendung.
\end{enumerate}

\subsection*{C: Berichtsrevision und Abschluss}

\begin{enumerate}
  \item Zunächst entstanden ein knapper Projektplan und zwei bearbeitbare
    Zwischenstände. Die Rückfragen des Verfassers zu Abbildungen, Bedienung und
    Formulierungen führten jeweils zu gezielten Überarbeitungen.
  \item Für den Abschluss wurden Berichtsvorlage, Protokoll und sämtliche
    Kommentare zum zweiten Zwischenstand ausgewertet. Daraus entstand die feste
    Gliederung; die Umsetzung jeder Rückmeldung ist in Anhang~C nachvollziehbar.
  \item Quellen zu Energy-Charts, PostgREST, Elm und verwandten Arbeiten wurden
    recherchiert und als BibTeX-Einträge eingebunden. Aussagen zu Einheiten und
    Datenbedeutung wurden auf Primärquellen zurückgeführt.
  \item Text, Tabellen und annotierte Screenshots wurden anschlie{\ss}end in LaTeX
    zusammengeführt. Der Bericht wurde kompiliert, vollständig gerendert und
    Seite für Seite visuell geprüft.
  \item Änderungen an Daten, Visualisierungen, Experimenten, Dokumentation und
    Bericht wurden in getrennten lokalen Git-Commits festgehalten. Ein Upload
    in das öffentliche Hochschul-GitLab erfolgt erst nach ausdrücklicher
    Freigabe.
\end{enumerate}

\section{Umsetzung des Feedbacks zum zweiten Zwischenstand}
\label{app:feedback}

\begin{longtable}{@{}p{0.32\textwidth}p{0.62\textwidth}@{}}
  \toprule
  Rückmeldung & Umsetzung \\
  \midrule
  \endhead
  Text ausbauen; Text direkt unter ``Einleitung'' & Zielproblem und
    Forschungsfrage stehen jetzt vor Unterabschnitt~1.1; alle Kapitel wurden
    fachlich ausgearbeitet. \\
  Forschungsfrage gehört in Kapitel~1 & In \cref{sec:einleitung} als
    hervorgehobene Leitfrage platziert. \\
  Anwendungshintergrund ergänzen & Unterschied zwischen physischem Fluss,
    Handel, Gebotszone, Erzeugungsmix und Preis in \cref{sec:hintergrund}
    erläutert. \\
  Einheiten nicht raten & Offizielle OpenAPI als Quelle; MW/GW-Transformation,
    Vorzeichen und EUR/MWh in \cref{sec:daten} dokumentiert. \\
  Zeitraum begründen & Vollständiger Mai 2025, Umfang und Kontraste in
    \cref{sec:daten} begründet. \\
  Anforderungen unabhängig von Techniken & Erst Aufgaben und Anforderungen in
    \cref{sec:aufgaben,sec:anforderungen}, danach Visualisierungsauswahl. \\
  Pfeilspitzen kleiner & SVG-Marker von 6 auf 4 Einheiten reduziert. \\
  maximale Linienbreite grö{\ss}er & Kodierung auf bis zu 21,5 SVG-Einheiten
    erweitert. \\
  Plus/Minus in Zeitreihe erklären & Sichtbare Vorzeichenlegende im
    Fluss-Teilplot und Erläuterung in \cref{sec:vis2}. \\
  andere Zeiträume auswählbar? & 48 Stunden, 7 Tage, Monat sowie Vor-/Zurück-
    Navigation umgesetzt. \\
  wei{\ss}e Matrixlücken prüfen & Zellkonturen entfernt; Nullwerte hellgrau statt
    wei{\ss} kodiert. \\
  Nutzen der Auswahl erklären & Matrix wählt Land und Stunde gleichzeitig;
    goldene Zeilen-/Spaltenmarkierung in \cref{sec:interaktion}. \\
  Farbtöne zeilenübergreifend vergleichbar? & Ein globales Monatsmaximum
    steuert alle Matrixzellen; numerische Legende ergänzt. \\
  Abbildung für jeden Anwendungsfall & Annotierte Abbildungen
    \cref{fig:case-france,fig:case-price,fig:case-matrix} ergänzt. \\
  \bottomrule
\end{longtable}
